\chapter*{Abstract}
This dissertation investigates knowledge distillation for efficient one-stage object detection using RetinaNet. Large detectors with \texttt{ResNet-50} or \texttt{ResNet-101} feature pyramid backbones typically provide stronger representations than lightweight models, but they are substantially more expensive to deploy under practical memory and latency constraints. The central aim of this work is therefore to determine whether a frozen teacher RetinaNet can transfer useful multi-scale structure to a smaller student detector, especially a \texttt{ResNet-18}-based RetinaNet, without increasing inference-time cost.

The proposed method combines the standard RetinaNet detection objective with a feature distillation term applied to aligned FPN representations. Rather than matching raw backbone activations, distillation is performed after the FPN so that teacher and student features already share a common semantic role and a common channel width. This allows the framework to compare selected pyramid levels directly without introducing an additional learned adaptation layer. Structural knowledge transfer is implemented through an SSIM-based loss on normalized, channel-aggregated feature maps, with the goal of preserving local spatial organization that is important for dense object detection.

At the current stage, the dissertation establishes the literature review, architectural design, dimensionality analysis, loss formulation, and COCO-style evaluation protocol needed for the experimental phase. The resulting framework provides a clear basis for comparing student-only RetinaNet training against teacher-guided training and for assessing whether hierarchy-aware structural distillation can improve compact detectors while keeping deployment efficient.
